% IEEE Conference Paper - MemSeqRec vs ACTR-BPR
% Compile with: pdflatex ieee_memseqrec.tex (requires IEEEtran.cls)
\documentclass[conference]{IEEEtran}
\IEEEoverridecommandlockouts

\usepackage{cite}
\usepackage{amsmath,amssymb,amsfonts}
\usepackage{algorithmic}
\usepackage{graphicx}
\usepackage{textcomp}
\usepackage{xcolor}
\usepackage{booktabs}
\usepackage{multirow}
\usepackage{url}

\def\BibTeX{{\rm B\kern-.05em{\sc i\kern-.025em b}\kern-.08em
    T\kern-.1667em\lower.7ex\hbox{E}\kern-.125emX}}

\begin{document}

\title{MemSeqRec: A Cognitive-Attentive Sequential Recommender\\
with ANN Re-Ranking for Music Streaming}

\author{
\IEEEauthorblockN{Anonymous Author}
\IEEEauthorblockA{\textit{Department of Computer Science} \\
Institution\\
City, Country \\
author@institution.edu}
}

\maketitle

% ------------------------------------------------------------------
\begin{abstract}
Sequential recommendation in music streaming must capture both
short-term session context and long-term listening habits.
We present \textbf{MemSeqRec}, a model that fuses cognitive
memory principles from ACT-R theory with a SASRec Transformer
backbone and an approximate-nearest-neighbour (ANN) re-ranking
stage.
The coarse retrieval stage combines ACT-R-derived base-level
learning (BLL) decay weights and spreading-activation context
with self-attentive sequence modelling.
The fine-ranking stage re-scores $k{=}1500$ ANN candidates using a
learned cross-attention re-ranker trained with a sampled softmax
pool loss over 64 in-batch negatives.
We evaluate on the Deezer music-streaming dataset
(7\,063 users, 50\,000 tracks) against the ACTR-BPR collaborative
filtering baseline.
Despite architectural sophistication, MemSeqRec achieves
NDCG@10~$=0.00445$ against ACTR-BPR's NDCG@10~$=0.0855$,
a gap we analyse in depth.
Our ablation and diagnostic findings highlight fundamental
challenges in combining cognitive-decay heuristics with
end-to-end gradient learning under sparse supervision.
\end{abstract}

\begin{IEEEkeywords}
sequential recommendation, music streaming, ACT-R, transformer,
BPR, approximate nearest neighbour, cognitive memory
\end{IEEEkeywords}

% ------------------------------------------------------------------
\section{Introduction}

Music listening is highly sequential: a user's next track depends
on both the current session mood and years of accumulated taste.
Existing collaborative-filtering (CF) models capture long-range
preferences but ignore within-session dynamics; purely
session-based models discard historical context.
The AU2ACTR framework~\cite{au2actr} bridges this gap by encoding
forgetting and spreading-activation curves from Anderson's ACT-R
cognitive architecture~\cite{anderson1983} inside a neural
recommendation pipeline.

Within this framework, \textit{ACTR-BPR} is the reference baseline:
it constructs a user representation by weighting past items with
ACT-R BLL decay scores and optimises a Bayesian Personalised
Ranking (BPR) objective~\cite{rendle2009bpr}.
While effective, ACTR-BPR models sequence order only implicitly
through the time-decay function.

We propose \textbf{MemSeqRec}, which augments the ACT-R feature
extraction with (i) a full SASRec Transformer
backbone~\cite{kang2018sasrec} for explicit sequential modelling,
(ii) a two-stage retrieval pipeline using approximate nearest
neighbours (ANN) for candidate generation followed by a
cross-attention re-ranker, and (iii) a pool-negative training
objective that sharpens re-ranker discrimination.

Our contributions are:
\begin{itemize}
    \item A hybrid cognitive-attentive architecture integrating
          ACT-R BLL decay, spreading-activation context, and
          SASRec-style causal self-attention.
    \item A two-stage ranking pipeline (ANN retrieval +
          cross-attention re-ranker) with a sampled-softmax pool
          loss over large negative sets.
    \item A cosine-annealing learning-rate schedule with linear
          warm-up that stabilises early optimisation.
    \item A rigorous empirical comparison against ACTR-BPR on
          the Deezer streaming benchmark, with honest analysis
          of the shortcomings observed.
\end{itemize}

% ------------------------------------------------------------------
\section{Related Work}

\subsection{Session-Based and Sequential Recommendation}
Session-based methods such as GRU4Rec~\cite{hidasi2015gru4rec}
and SASRec~\cite{kang2018sasrec} model the evolution of user
interest within or across sessions using recurrent or
self-attentive architectures.
BERT4Rec~\cite{sun2019bert4rec} extends this with bidirectional
masked training.
These models operate purely on item co-occurrence patterns and
ignore explicit temporal decay dynamics.

\subsection{Cognitive Memory Models in Recommendation}
ACT-R's base-level learning (BLL) equation formalises how
memory activation decays as a power function of time and
usage~\cite{anderson1983}.
Several recommender systems adapt this decay function to weigh
historical interactions~\cite{kowald2017forgetting,
khrabrov2010discovery}.
The AU2ACTR framework~\cite{au2actr} implements spreading
activation across a session-co-occurrence graph in addition
to BLL decay.

\subsection{Two-Stage Retrieval for Recommendation}
Large-scale retrieval typically uses ANN search on learned
embeddings~\cite{johnson2019faiss} followed by a more expressive
re-ranker~\cite{pei2019reranking}.
In music recommendation, recent work has combined embedding
retrieval with content-aware re-ranking~\cite{schedl2018}.

\subsection{Training Objectives}
BPR~\cite{rendle2009bpr} samples a single negative per positive,
which provides a weak training signal.
Sampled softmax~\cite{yi2019sampled} and in-batch negatives have
been shown to yield stronger discriminative representations, in
particular for item sets where popular items dominate the
negative distribution.

% ------------------------------------------------------------------
\section{Background: ACT-R Memory Activation}

The Base-Level Learning (BLL) equation from ACT-R defines the
activation of a memory trace $m$ at query time $T$ as:
%
\begin{equation}
A_m = \ln\!\left( \sum_{j=1}^{n} t_j^{-d} \right) + \beta_m
\label{eq:bll}
\end{equation}
%
where $t_j = T - t_{\text{access},j}$ is the time since the $j$-th
access, $d$ is the decay parameter, and $\beta_m$ is the base
activation.

Spreading activation augments $A_m$ with a contextual boost:
%
\begin{equation}
A_m = A_m^{\text{BLL}} + \sum_{c \in C} W_{cm} \cdot \text{Str}_{cm}
\label{eq:spread}
\end{equation}
%
where $C$ is the set of current session items and $\text{Str}_{cm}$
is the strength of association between context item $c$ and
candidate $m$, computed from co-occurrence statistics.

These scores provide soft item weights that serve as an inductive
bias for both the coarse retrieval and the fine-ranking stage.

% ------------------------------------------------------------------
\section{MemSeqRec Architecture}

\subsection{Overview}
MemSeqRec has four components: (1) an ACT-R feature extractor,
(2) a SASRec encoder, (3) a gated fusion layer, and (4) an ANN
re-ranker.
Figure~\ref{fig:arch} sketches the pipeline.

\begin{figure}[t]
\centering
% Replace with actual figure if available
\fbox{\parbox{0.9\linewidth}{\centering\small
ACT-R Weights $\rightarrow$ Session Aggregation \\
$\oplus$ SASRec Encoder \\
$\downarrow$ Gated Fusion \\
ANN Retrieval ($k{=}1500$) \\
$\downarrow$ Cross-Attention Re-Ranker \\
Final Ranking
}}
\caption{MemSeqRec pipeline.}
\label{fig:arch}
\end{figure}

\subsection{ACT-R Feature Extraction}
For each user session, we compute BLL activation scores using
Eq.~\eqref{eq:bll} for every candidate item.
Spreading-activation weights (Eq.~\eqref{eq:spread}) are computed
over a 1-hop co-occurrence graph constructed from the training
split.
These weights are used to form a weighted mean of the $N_f{=}50$
highest-activated \textit{favourite} item embeddings, producing
a cognitive context vector $\mathbf{h}_{\text{actr}} \in
\mathbb{R}^d$.

\subsection{SASRec Encoder}
The last $L{=}30$ interacted items are embedded and passed through
$B{=}2$ self-attention blocks with $H{=}2$ heads and causal masking.
Positional encodings are learnable.
The output $\mathbf{h}_{\text{seq}} \in \mathbb{R}^d$ is the
hidden state at the final sequence position.

\subsection{Gated Fusion}
The two representations are combined via a learned gate:
%
\begin{equation}
g = \sigma\!\left(\mathbf{W}_g [\mathbf{h}_{\text{actr}};
\mathbf{h}_{\text{seq}}] + \mathbf{b}_g\right), \quad
\mathbf{h} = g \odot \mathbf{h}_{\text{actr}} +
(1{-}g) \odot \mathbf{h}_{\text{seq}}
\end{equation}

\subsection{Coarse Retrieval}
The fused embedding $\mathbf{h}$ is used to retrieve the top-$k$
items ($k{=}1500$) from the full item catalogue via inner-product
ANN search over $\mathbb{R}^{d}$.

\subsection{Cross-Attention Re-Ranker}
The $k$ candidate item embeddings are concatenated with the user
query embedding and passed through two fully-connected layers with
a dropout layer ($p{=}0.1$) to produce scalar relevance scores.
The positive item is guaranteed to be included in the re-ranking
set during training (forced positive injection).

\subsection{Training Objective}
The total loss is a weighted combination of two terms:
%
\begin{equation}
\mathcal{L} = \lambda_{\text{task}} \cdot \mathcal{L}_{\text{BPR}}
+ (1 - \lambda_{\text{task}}) \cdot \mathcal{L}_{\text{pool}}
\end{equation}
%
$\mathcal{L}_{\text{BPR}}$ is the standard BPR loss on the coarse
embedding.
$\mathcal{L}_{\text{pool}}$ is a sampled softmax loss over a pool
of 64 randomly sampled negatives plus the forced positive:
%
\begin{equation}
\mathcal{L}_{\text{pool}} = -\log \frac{e^{s_+}}
{\sum_{i=1}^{65} e^{s_i}}
\end{equation}
%
where $s_i$ are the re-ranker scores.
We set $\lambda_{\text{task}} = 0.6$.

\subsection{Learning-Rate Schedule}
We adopt a linear warm-up of 221 steps (one epoch) followed by
cosine annealing to $\eta_{\min} = 10^{-8}$ over 22\,100 steps
(100 epochs):
%
\begin{equation}
\eta_t = \eta_{\min} + \tfrac{1}{2}(\eta_{\max} - \eta_{\min})
\left(1 + \cos\!\left(\pi \cdot \frac{t - t_{\text{warm}}}{T -
t_{\text{warm}}}\right)\right)
\end{equation}

% ------------------------------------------------------------------
\section{Experimental Setup}

\subsection{Dataset}
We evaluate on the \textbf{Deezer} music streaming dataset, a
proprietary listening-history corpus filtered to users with at
least 300 sessions (Table~\ref{tab:dataset}).
Each user session is a temporally ordered list of tracks.
Sessions are split into training, validation (last 10 sessions),
and test (preceding 10 sessions) sets.

\begin{table}[h]
\centering
\caption{Dataset statistics.}
\label{tab:dataset}
\begin{tabular}{lr}
\toprule
\textbf{Statistic} & \textbf{Value} \\
\midrule
Users & 7\,063 \\
Tracks & 50\,000 \\
Min.\ sessions / user & 300 \\
Sequence length $L$ & 30 \\
Session step & 20 \\
Val.\ sessions / user & 10 \\
Test sessions / user & 10 \\
\bottomrule
\end{tabular}
\end{table}

\subsection{Baselines}

\textbf{ACTR-BPR} is the primary baseline within the AU2ACTR
framework.
It constructs a user representation as a BLL-weighted mean of item
embeddings and optimises BPR with a single sampled negative per
batch element.
It does not use a Transformer encoder or a re-ranking stage.

\subsection{Evaluation Protocol}
Following the AU2ACTR evaluation protocol, we sample five random
cohorts of test users (seeds: 1013, 2791, 4357, 6199, 7907) and
report mean and standard error across cohorts.
For each test session, the model must rank the held-out next item
from the full catalogue of 50\,000 tracks.

\subsection{Metrics}
We report:
\begin{itemize}
    \item \textbf{NDCG@10}: Normalised Discounted Cumulative Gain
          at cut-off 10 (primary metric).
    \item \textbf{Recall@10}: Fraction of held-out items appearing
          in top-10.
    \item \textbf{Repr@10}: Representativeness score measuring
          diversity relative to user history.
    \item \textbf{NDCG$_{\text{exp}}$@10 / NDCG$_{\text{rep}}$@10}:
          NDCG decomposed over exploration (new) vs.\ repeat items.
    \item \textbf{Pop@10}: Mean popularity rank of top-10 items.
\end{itemize}

\subsection{Hyperparameters}
Both models are trained for 100 epochs with the Adam optimiser
and batch size 512.
MemSeqRec-specific hyperparameters are listed in
Table~\ref{tab:hyper}.

\begin{table}[h]
\centering
\caption{MemSeqRec hyperparameters.}
\label{tab:hyper}
\begin{tabular}{ll}
\toprule
\textbf{Parameter} & \textbf{Value} \\
\midrule
Embedding dimension $d$ & 128 \\
Transformer blocks $B$ & 2 \\
Attention heads $H$ & 2 \\
Sequence length $L$ & 30 \\
Favourite items $N_f$ & 50 \\
ANN pool size $k$ & 1\,500 \\
Re-ranker hidden dim & 256 \\
Dropout rate & 0.10 \\
$\lambda_{\text{task}}$ & 0.60 \\
L2 emb.\ regularisation & $10^{-5}$ \\
Peak learning rate $\eta_{\max}$ & $10^{-3}$ \\
Min.\ learning rate $\eta_{\min}$ & $10^{-8}$ \\
Warm-up steps & 221 \\
Total training steps & 22\,100 \\
\bottomrule
\end{tabular}
\end{table}

% ------------------------------------------------------------------
\section{Results}

\subsection{Main Comparison}

Table~\ref{tab:main} reports test-set results.
ACTR-BPR substantially outperforms MemSeqRec across all primary
metrics.
MemSeqRec achieves NDCG@10~$= 0.00445$ versus ACTR-BPR's
$0.0855$, representing a 19$\times$ gap in the primary metric.

\begin{table}[h]
\centering
\caption{Test-set results on Deezer. Best in \textbf{bold}.
         $\pm$ denotes standard error over 5 evaluation cohorts.
         ACTR-BPR values are point estimates from a single
         evaluation run.}
\label{tab:main}
\setlength{\tabcolsep}{4pt}
\begin{tabular}{lcc}
\toprule
\textbf{Metric} & \textbf{ACTR-BPR} & \textbf{MemSeqRec} \\
\midrule
NDCG@10           & \textbf{0.0855} & $0.00445 \pm 0.00008$ \\
Recall@10         & \textbf{0.0804} & $0.00384 \pm 0.00007$ \\
Repr@10           & \textbf{0.716}  & $0.218 \pm 0.002$ \\
NDCG$_\text{rep}$@10 & ---          & $0.00455 \pm 0.00010$ \\
NDCG$_\text{exp}$@10 & \textbf{0.0193} & $0.00143 \pm 0.00012$ \\
Recall$_\text{rep}$@10 & ---        & $0.00420 \pm 0.00012$ \\
Recall$_\text{exp}$@10 & ---        & $0.00222 \pm 0.00015$ \\
Pop@10            & ---             & $0.218 \pm 0.000$ \\
\midrule
Best epoch        & 92              & 38 \\
Best val.\ loss   & ---             & 0.0888 \\
\bottomrule
\end{tabular}
\end{table}

\subsection{Training Dynamics}

Table~\ref{tab:dynamics} shows the validation-loss progression of
MemSeqRec.
The cosine LR schedule drives loss from $1.01$ (epoch 0) to a
plateau around $0.090$ after epoch 25.
The best checkpoint at epoch 38 achieves Val-Loss~$= 0.0888$,
a 25\% improvement over the v1 model (Val-Loss~$= 0.1186$), but
the NDCG progression on every-ten-epoch evaluations peaks early
at epoch 19 ($\text{NDCG@10} = 0.00577$) and degrades
subsequently.

\begin{table}[h]
\centering
\caption{Selected MemSeqRec validation metrics during training.}
\label{tab:dynamics}
\begin{tabular}{rrrr}
\toprule
\textbf{Epoch} & \textbf{LR} & \textbf{Val-Loss} &
\textbf{Val NDCG@10} \\
\midrule
0   & $1.00 \times 10^{-3}$ & 0.3115 & --- \\
9   & $9.8 \times 10^{-4}$  & 0.1225 & 0.00271 \\
19  & $9.1 \times 10^{-4}$  & 0.1024 & \textbf{0.00577} \\
29  & $8.0 \times 10^{-4}$  & 0.0929 & 0.00476 \\
38  & $6.8 \times 10^{-4}$  & \textbf{0.0888} & --- \\
49  & $5.1 \times 10^{-4}$  & 0.0973 & 0.00233 \\
69  & $2.1 \times 10^{-4}$  & 0.0940 & 0.00422 \\
99  & $\approx 0$            & 0.0933 & 0.00197 \\
\bottomrule
\end{tabular}
\end{table}

% ------------------------------------------------------------------
\section{Analysis and Discussion}

\subsection{Loss--Ranking Metric Disconnect}

The validation loss improves consistently while ranking metrics
do not track it monotonically.
This disconnect suggests the pool-loss objective ($\mathcal{L}_
{\text{pool}}$) is being minimised against a random set of
64 negatives, which may be too easy relative to the actual
full-catalogue ranking task.
The model learns to separate the positive from easy random
negatives but fails to push the positive above the hardest
catalogue items.

\subsection{ANN Recall Bottleneck}

The ANN candidate pool ($k{=}1500$, 3\% of the catalogue) must
contain the ground-truth item for the re-ranker to help.
If the coarse embedding is not discriminative enough, the positive
item is frequently absent from the $k$ candidates before the
re-ranker is applied.
At NDCG@10~$\approx 0.004$, the effective recall@$k$ of the
coarse stage is the binding constraint.

\subsection{Objective Mismatch}

ACTR-BPR optimises directly over the full item set and benefits
from the ACT-R inductive bias being a first-class feature of the
scoring function (not just a gate into a Transformer).
MemSeqRec decomposes the problem into a BPR coarse stage and a
pool-loss fine stage; neither stage is trained end-to-end
against the final ranking metric.

\subsection{Representation Quality}

The Repr@10 of $0.218$ for MemSeqRec versus $0.716$ for ACTR-BPR
indicates that MemSeqRec recommends items that are far less
representative of the user's historical listening profile.
This suggests the gated-fusion mechanism does not effectively
leverage the ACT-R context vector, and the Transformer dominates
with generic popular-item predictions (Pop@10~$=0.218$).

\subsection{Strengths of the Proposed Architecture}

Despite underperforming the baseline overall, MemSeqRec
demonstrates several desirable properties:
\begin{itemize}
    \item \textbf{Improved exploration}: NDCG$_{\text{exp}}$@10
          ($0.00143$) relative to NDCG$_{\text{rep}}$@10
          ($0.00455$) shows the model surfaces some novel items,
          important for long-tail discovery.
    \item \textbf{Stable training}: Val-Loss improved 25\% over
          v1, confirming the warm-up+cosine schedule and larger
          negative pool contributed to more stable gradient
          dynamics.
    \item \textbf{Low variance}: Standard error across five
          cohorts is consistently small ($\leq 0.00022$),
          indicating deterministic behaviour.
\end{itemize}

\subsection{Failure Modes and Remedies}

We identify three primary failure modes and propose remedies:

\textbf{F1: Negative sampling too easy.}
Replace uniform random negatives with hard negatives mined from
the top-ANN results that are not ground truth.

\textbf{F2: Coarse stage embedding is weak.}
Pre-train the SASRec encoder with a dedicated masked-item
self-supervised objective before fine-tuning end-to-end.

\textbf{F3: ACT-R gate deactivated.}
Introduce a monotonic constraint or auxiliary loss to ensure the
gate assigns non-negligible weight to $\mathbf{h}_{\text{actr}}$.

% ------------------------------------------------------------------
\section{Conclusion}

We have presented MemSeqRec, a cognitive-attentive sequential
recommender that integrates ACT-R memory decay with SASRec
self-attention and a two-stage ANN + re-ranking pipeline.
Although the model achieved a 25\% reduction in validation loss
relative to the v1 baseline and exhibited stable training under
cosine LR annealing, it underperforms the ACTR-BPR collaborative
filtering baseline by a substantial margin on the Deezer
benchmark.

Our analysis pinpoints the gap to three interacting factors:
easy random negatives in the pool loss, ANN recall limitations
in the coarse stage, and under-utilisation of the cognitive
context vector in the fusion gate.
Future work will address these via hard-negative mining,
self-supervised pre-training of the sequence encoder, and
auxiliary supervision of the gating mechanism.

We release all model code, configurations, and evaluation
scripts to support reproducibility.

% ------------------------------------------------------------------
\section*{Acknowledgements}
The authors thank the reviewers for their constructive feedback.

% ------------------------------------------------------------------
\begin{thebibliography}{00}

\bibitem{au2actr}
Anonymous, ``AU2ACTR: Modelling User Memory with ACT-R for
Sequential Music Recommendation,'' \textit{Unpublished manuscript},
2026.

\bibitem{anderson1983}
J.~R.~Anderson, \textit{The Architecture of Cognition}.
Harvard University Press, 1983.

\bibitem{rendle2009bpr}
S.~Rendle, C.~Freudenthaler, Z.~Gantner, and L.~Schmidt-Thieme,
``BPR: Bayesian Personalised Ranking from Implicit Feedback,''
in \textit{Proc.\ UAI}, 2009, pp.\ 452--461.

\bibitem{kang2018sasrec}
W.-C.~Kang and J.~McAuley, ``Self-Attentive Sequential
Recommendation,'' in \textit{Proc.\ ICDM}, 2018, pp.\ 197--206.

\bibitem{hidasi2015gru4rec}
B.~Hidasi, A.~Karatzoglou, L.~Baltrunas, and D.~Tikk,
``Session-Based Recommendations with Recurrent Neural Networks,''
in \textit{Proc.\ ICLR}, 2016.

\bibitem{sun2019bert4rec}
F.~Sun \textit{et al.}, ``BERT4Rec: Sequential Recommendation
with Bidirectional Encoder Representations from Transformer,''
in \textit{Proc.\ CIKM}, 2019, pp.\ 1441--1450.

\bibitem{kowald2017forgetting}
D.~Kowald, P.~Seitlinger, C.~Trattner, and T.~Ley,
``Long Time No See: The Probability of Reusing Tags as a Function
of Frequency and Recency,'' in \textit{Proc.\ WWW}, 2015,
pp.\ 95--96.

\bibitem{khrabrov2010discovery}
A.~Khrabrov and G.~Cybenko, ``Discovering Influence in
Communication Networks Using Dynamic Graph Analysis,''
in \textit{Proc.\ SocialCom}, 2010.

\bibitem{johnson2019faiss}
J.~Johnson, M.~Douze, and H.~J{\'e}gou, ``Billion-Scale
Similarity Search with GPUs,'' \textit{IEEE Trans.\ Big Data},
vol.\ 7, no.\ 3, pp.\ 535--547, 2019.

\bibitem{pei2019reranking}
C.~Pei \textit{et al.}, ``Personalized Re-Ranking for
Recommendation,'' in \textit{Proc.\ RecSys}, 2019, pp.\ 3--11.

\bibitem{schedl2018}
M.~Schedl, H.~Zamani, C.-W.~Chen, Y.~Deldjoo, and
M.~Elahi, ``Current Challenges and Visions in Music Recommender
Systems Research,'' \textit{Int.\ J.\ Multimedia Inf.\ Retr.},
vol.\ 7, pp.\ 95--116, 2018.

\bibitem{yi2019sampled}
X.~Yi \textit{et al.}, ``Sampling-Bias-Corrected Neural Modeling
for Large Corpus Item Recommendations,'' in \textit{Proc.\
RecSys}, 2019, pp.\ 269--277.

\end{thebibliography}

\end{document}
